\chapter{Erythrocyte Sedimentation Rate (ESR)}

\section*{What Is This Test?}
The erythrocyte sedimentation rate (ESR), also known as the sed rate, measures the rate at which red blood cells (erythrocytes) settle in a vertical column of anticoagulated whole blood over a specified period, typically one hour. It is a nonspecific indicator of inflammation and is one of the oldest laboratory tests still in widespread clinical use. The ESR reflects changes in plasma protein composition, specifically elevations in fibrinogen and other acute-phase reactant proteins (alpha-1-antitrypsin, haptoglobin, C-reactive protein, ceruloplasmin) that reduce the negative charge (zeta potential) between red cells, promoting rouleaux formation --- the stacking of red blood cells like coins. Rouleaux are heavier aggregates that sediment more rapidly under gravity. Immunoglobulins (IgG, IgM), particularly in monoclonal gammopathies such as multiple myeloma and Waldenstr\"{o}m macroglobulinemia, also strongly promote rouleaux and markedly elevated ESR. The ESR is used primarily to detect and monitor inflammatory conditions, including giant cell (temporal) arteritis, polymyalgia rheumatica, rheumatoid arthritis, systemic lupus erythematosus, chronic infections (osteomyelitis, endocarditis, tuberculosis), and certain malignancies. It is particularly valuable in the diagnosis of temporal arteritis and polymyalgia rheumatica, where a markedly elevated ESR is a diagnostic criterion. However, because the ESR is nonspecific (elevated in many conditions), it must be interpreted in the context of the clinical presentation and other laboratory tests, most notably C-reactive protein (CRP) \cite{brigden1999clinical}.

\section*{Alternative Names}
ESR; Sed Rate; Erythrocyte Sedimentation Rate; Westergren Sedimentation Rate; Westergren ESR

\section*{Principle}
The ESR is measured by the Westergren method, the international reference method established by the International Council for Standardization in Haematology (ICSH) \cite{icsh2011icsh}. Undiluted whole blood is collected in EDTA (lavender-top tube) or sodium citrate (light blue-top tube, 3.2\% trisodium citrate in a 1:4 citrate-to-blood ratio). For the traditional Westergren method, the anticoagulated blood is diluted 1:4 with 3.8\% trisodium citrate (1 part citrate, 4 parts blood), and a Westergren pipette (a calibrated glass or plastic tube 200 mm long with an internal diameter of 2.55 mm) is filled by suction to the zero mark. The pipette is placed in a vertical position in a leveling rack and left undisturbed at room temperature for exactly 60 minutes. The distance (in millimeters) from the top of the column of plasma to the top of the sedimented red cell layer is measured visually and reported as mm/hour. For the modified Westergren method, EDTA-anticoagulated blood is used directly, diluted 1:4 with citrate or saline, and read at 60 minutes. Automated ESR analyzers (StaRRsed, Ves-Matic, Sediplast) use similar sedimentation principles but with automated reading and shorter analysis times (20--30 minutes) that are mathematically converted to 1-hour Westergren equivalents. The ESR is influenced by hematocrit (anemia causes a falsely elevated ESR because fewer red cells encounter less frictional resistance during sedimentation; polycythemia causes a falsely low ESR), plasma proteins (fibrinogen, immunoglobulins, and acute-phase reactants promote rouleaux), and red blood cell morphology (sickle cells, spherocytes, and marked anisopoikilocytosis interfere with rouleaux formation and decrease ESR). The ESR does not directly measure any single molecule; it is a composite reflection of the balance between pro-sedimentation factors (fibrinogen and large asymmetric proteins) and factors that resist sedimentation (the negative charge on red cell membranes, or zeta potential) \cite{tietz2023textbook}.

\section*{History}
The erythrocyte sedimentation rate is one of the oldest clinical laboratory tests, first described by the Polish physician Edmund Faustyn Biernacki in 1897. Biernacki noted that red blood cells settled at different rates in different individuals and that the rate was altered in disease states. The phenomenon was independently rediscovered and popularized by the Swedish hematologist Robert Sanno F{\aa}hraeus in 1918, who demonstrated that the sedimentation rate was elevated in pregnancy and in various diseases, and attributed the increase to elevated fibrinogen. The standardized Westergren method was developed by Alf Vilhelm Westergren in 1921, using sodium citrate as anticoagulant and a 200 mm vertical glass tube; his method became the international standard. In 1926, Westergren published definitive clinical data establishing the normal values and clinical utility of the ESR. Throughout the 20th century, the ESR was a standard screening test for occult disease. While its clinical role has been partially eclipsed by C-reactive protein (CRP) --- which is more sensitive, rises and falls more rapidly, and is not affected by hematocrit or red cell morphology --- the ESR remains a valuable test, particularly for the diagnosis and monitoring of temporal arteritis, polymyalgia rheumatica, and certain chronic inflammatory conditions where ESR and CRP are not always concordant.

\section*{Why This Test Is Ordered}
\begin{itemize}
  \item Diagnosis of giant cell (temporal) arteritis --- markedly elevated ESR (often >50 mm/hour, frequently >100 mm/hour) is a key diagnostic criterion, and the ESR is used to monitor response to high-dose corticosteroid therapy. A normal ESR makes temporal arteritis unlikely but does not exclude it \cite{salvarani2008polymyalgia}.
  \item Diagnosis of polymyalgia rheumatica --- ESR is frequently >40 mm/hour and is a diagnostic criterion; used to monitor disease activity and response to corticosteroid therapy \cite{kermani2013polymyalgia}
  \item Screening for and monitoring of systemic inflammatory and autoimmune diseases --- rheumatoid arthritis (ESR often elevated and correlates with disease activity), systemic lupus erythematosus (ESR elevated but may not correlate with clinical activity as well as CRP does), ankylosing spondylitis, psoriatic arthritis, and inflammatory bowel disease
  \item Diagnosis and monitoring of chronic infections --- osteomyelitis, subacute bacterial endocarditis, tuberculosis, brucellosis, and deep-seated abscess
  \item Evaluation of fever of unknown origin (FUO) --- an elevated ESR supports an inflammatory, infectious, or neoplastic cause
  \item Monitoring of certain malignancies --- Hodgkin lymphoma and non-Hodgkin lymphoma, renal cell carcinoma, multiple myeloma (markedly elevated ESR with rouleaux on peripheral smear)
  \item Diagnosis of multiple myeloma and Waldenstr\"{o}m macroglobulinemia --- markedly elevated ESR (often >100 mm/hour) due to paraprotein-induced rouleaux formation
  \item Assessment of disease activity and response to therapy in chronic inflammatory conditions
  \item Differentiation of inflammatory from non-inflammatory conditions --- useful in the initial evaluation of nonspecific symptoms such as weight loss, fatigue, and generalized malaise
  \item Part of the diagnostic workup for suspected osteomyelitis, especially in diabetic foot infections and vertebral osteomyelitis
\end{itemize}

\section*{Is This a Primary or Secondary Test}
The ESR is a nonspecific secondary test that serves as a marker of inflammation and is used to support the diagnosis and monitor the activity of specific inflammatory, infectious, and neoplastic conditions.

\section*{How This Test Is Performed}
A venous blood sample is collected in a lavender-top (EDTA) tube from a peripheral vein. Approximately 2--5 mL of blood is drawn, and the tube is gently inverted 8--10 times to ensure adequate mixing. For the Westergren method, the blood is diluted with either 3.8\% sodium citrate (1:4 ratio) or 0.9\% saline. A Westergren pipette or tube (calibrated 200 mm column) is filled to the zero mark. The tube is placed vertically in a rack and left undisturbed at room temperature for exactly 60 minutes. After 60 minutes, the distance from the top of the plasma column to the top of the sedimented red blood cell layer is measured in millimeters and reported as mm/hour. Automated ESR analyzers use a similar sedimentation principle but employ infrared sensors or digital imaging to track the sedimentation interface. Some automated instruments use undiluted EDTA blood and read at shorter intervals (20--30 minutes), converting the result mathematically to the standard 1-hour Westergren equivalent \cite{fischbach2023manual}. The venipuncture takes less than 5 minutes; testing requires 1 hour for the manual Westergren method or 20--30 minutes for automated methods.

\section*{What Preparation Is Needed}
No special preparation is required. Fasting is not necessary for the ESR. However, if other tests ordered concurrently require fasting (fasting glucose, lipid profile), the patient should follow those fasting requirements. A high-fat meal just before testing may cause lipemia, which can affect the ESR (mild elevation). Patients should inform their healthcare provider of all medications, particularly corticosteroids (which lower ESR), oral contraceptives and estrogen therapy (which increase ESR), aspirin and NSAIDs (which may decrease ESR in inflammatory conditions), and biologic therapies for autoimmune disease. Recent immunization, surgery, or trauma can transiently elevate the ESR.

\section*{Contraindications and Precautions}
No absolute contraindications. Numerous pre-analytical and patient-related factors affect the ESR and must be considered for accurate interpretation: the sample must be thoroughly mixed before testing; testing should be performed within 4 hours of collection at room temperature or 24 hours if refrigerated at 2--8\textdegree{}C (with re-warming to room temperature before testing); hemolyzed, clotted, or insufficiently filled specimens are unacceptable; the Westergren tube must be perfectly vertical --- a tilt of as little as 3 degrees can increase the ESR by up to 30\%; vibration (nearby centrifuges, heavy machinery) can artificially increase the ESR; room temperature extremes --- high temperatures (>25\textdegree{}C) increase ESR, low temperatures (<18\textdegree{}C) decrease ESR; anemia (hematocrit <35\%) causes a falsely elevated ESR because fewer red blood cells settle more rapidly; polycythemia (hematocrit >50\%) causes a falsely low ESR; red blood cell morphology abnormalities --- sickle cells (which cannot form rouleaux due to their rigid, elongated shape) produce a markedly low ESR even in the presence of severe inflammation; spherocytes, acanthocytes, and marked poikilocytosis also interfere with rouleaux and decrease ESR; extremely elevated fibrinogen or immunoglobulins can be beyond the linear range of the ESR; ESR is nonspecific and does not identify the cause of inflammation --- it should not be used for general health screening in asymptomatic individuals.

\section*{Risks}
Minimal risks of venipuncture: mild pain or stinging at the needle site, bruising, hematoma, lightheadedness or vasovagal syncope, and extremely rarely, infection or nerve injury.

\section*{What Happens After the Test}
Results from automated analyzers are available within 30--60 minutes; manual Westergren results require exactly 1 hour of sedimentation time plus processing. The ESR result is reported in mm/hour. The healthcare provider interprets the ESR in the context of clinical symptoms, physical examination findings, and other laboratory tests, particularly the C-reactive protein (CRP). An elevated ESR supports the presence of an inflammatory, infectious, or neoplastic process but does not identify the specific cause. In the diagnosis of temporal arteritis, a markedly elevated ESR (>50 mm/hour, often >100 mm/hour) combined with compatible clinical symptoms (new headache, scalp tenderness, jaw claudication, visual disturbance in a patient >50 years) supports the diagnosis, and temporal artery biopsy should be performed. High-dose corticosteroids should be started immediately if temporal arteritis is suspected, without waiting for biopsy results \cite{salvarani2008polymyalgia}. Serial ESR measurements are used to monitor response to therapy and detect disease relapse \cite{tierney2023current}.

\section*{Understanding Results}

\subsection*{Below Normal (Decreased ESR)}
\begin{itemize}
  \item \textbf{Polycythemia:} Increased red blood cell mass (polycythemia vera, secondary polycythemia) increases the frictional resistance between settling red cells, decreasing the ESR.
  \item \textbf{Sickle cell disease:} Rigid, elongated sickle-shaped red blood cells cannot form rouleaux, producing a profoundly low ESR (often 0--5 mm/hour) even in the presence of severe infection, osteomyelitis, or vaso-occlusive crisis. A normal or elevated ESR in a patient with sickle cell disease may indicate a disease complication.
  \item \textbf{Spherocytosis (hereditary spherocytosis, autoimmune hemolytic anemia):} Spherical red blood cells have decreased surface area-to-volume ratio, cannot form rouleaux effectively, and sediment poorly.
  \item \textbf{Hypofibrinogenemia or afibrinogenemia:} Congenital or acquired (DIC, liver disease) fibrinogen deficiency eliminates the primary protein promoting rouleaux formation.
  \item \textbf{Polycythemia of the newborn:} Physiologically high hematocrit in neonates decreases the ESR.
  \item \textbf{Corticosteroid therapy:} Glucocorticoids suppress the acute-phase response and decrease fibrinogen and other pro-sedimentation proteins.
  \item \textbf{Extreme leukocytosis:} Very high white blood cell counts (>100,000/$\mu$L) in leukemia may pack into the upper plasma layer and inhibit red cell sedimentation.
  \item \textbf{Cryoglobulinemia:} Cryoprecipitate proteins at room temperature may interfere with the plasma-column interface.
  \item \textbf{Technically invalid:} Clotted sample, tilted Westergren tube, expired blood sample (>4 hours unrefrigerated).
\end{itemize}

\subsection*{Above Normal (Increased ESR)}
\begin{itemize}
  \item \textbf{Giant cell (temporal) arteritis:} The cardinal laboratory abnormality; ESR is often >50 mm/hour and frequently >100 mm/hour. Nearly 90\% of patients have ESR >50 mm/hour. A normal ESR does not exclude temporal arteritis (10--20\% of biopsy-proven cases have normal ESR). The ESR is used to monitor disease activity and response to corticosteroid therapy.
  \item \textbf{Polymyalgia rheumatica:} ESR is frequently >40 mm/hour and often >80 mm/hour. Isolated proximal muscle pain and stiffness (shoulders, hips, neck) in an older adult (>50 years) with markedly elevated ESR strongly supports the diagnosis. Dramatic and rapid improvement with low-dose corticosteroids (prednisone 10--20 mg/day) is characteristic.
  \item \textbf{Multiple myeloma and Waldenstr\"{o}m macroglobulinemia:} Monoclonal immunoglobulins (IgG, IgA, IgM) strongly promote rouleaux formation. ESR is often >100 mm/hour and may be the presenting laboratory abnormality in otherwise asymptomatic patients. Rouleaux are prominent on peripheral blood smear.
  \item \textbf{Rheumatoid arthritis:} ESR correlates with disease activity and is used with joint counts, patient global assessment, and CRP to calculate disease activity scores (DAS28-ESR) \cite{sokka2009erythrocyte}. Guiding therapy to achieve clinical remission often involves normalizing or substantially reducing the ESR.
  \item \textbf{Systemic lupus erythematosus:} ESR is often elevated, especially during active disease, but may not correlate as well with clinical activity as CRP and complement levels (C3, C4).
  \item \textbf{Chronic infections:} Osteomyelitis (especially vertebral osteomyelitis), subacute bacterial endocarditis, tuberculosis, brucellosis, deep-seated abscesses, and fungal infections.
  \item \textbf{Inflammatory bowel disease:} Crohn disease and ulcerative colitis; ESR correlates with disease activity and response to therapy.
  \item \textbf{Malignancy:} Hodgkin lymphoma (ESR is a prognostic factor in early-stage disease), non-Hodgkin lymphoma, renal cell carcinoma (paraneoplastic), and any metastatic malignancy.
  \item \textbf{Pregnancy:} Physiological elevation, especially in the third trimester; ESR may reach 30--50 mm/hour.
  \item \textbf{Anemia:} Any cause of anemia (iron deficiency, B12/folate deficiency, hemolysis, chronic disease) decreases hematocrit and artificially elevates ESR. The ESR result should be corrected for hematocrit or interpreted with this known effect in mind.
  \item \textbf{Advanced age:} ESR increases mildly with age; the formula for the upper limit of normal in older adults is: Men: Age / 2; Women: (Age + 10) / 2.
  \item \textbf{Acute myocardial infarction:} ESR rises 2--3 days post-MI as part of the acute-phase response and may remain elevated for weeks.
  \item \textbf{Postoperative state and tissue injury:} ESR rises within 2--3 days after surgery or trauma and normalizes over several weeks.
  \item \textbf{Chronic kidney disease and nephrotic syndrome:} Hypoalbuminemia and dysproteinemia alter rouleaux.
  \item \textbf{Hypothyroidism and hyperthyroidism:} Thyroid dysfunction can alter protein metabolism and ESR.
\end{itemize}

\section*{Normal Results}
\begin{itemize}
  \item \textbf{Adult males (<50 years):} 0--15 mm/hour
  \item \textbf{Adult females (<50 years):} 0--20 mm/hour
  \item \textbf{Adult males (>50 years):} 0--20 mm/hour or calculated as Age/2
  \item \textbf{Adult females (>50 years):} 0--30 mm/hour or calculated as (Age + 10)/2
  \item \textbf{Children (1--12 years):} 0--10 mm/hour
  \item \textbf{Neonates:} 0--2 mm/hour (very low due to polycythemia of the newborn)
  \item \textbf{Pregnant women:} First trimester 0--25 mm/hour; second trimester 0--45 mm/hour; third trimester 0--50 mm/hour
  \item \textbf{Note:} The ESR increases with age in healthy individuals, and age-adjusted normal values should be used for older patients \cite{wallach2022interpretation}.
\end{itemize}

\section*{Follow-up Tests}
\begin{itemize}
  \item \textbf{C-reactive protein (CRP):} A more sensitive, specific, and rapidly changing marker of inflammation. CRP rises within 4--6 hours of inflammatory stimulus and falls within 24--48 hours of its resolution (unlike ESR, which takes days to weeks). CRP is not affected by hematocrit, red cell morphology, or immunoglobulins. ESR and CRP are complementary and should both be ordered when temporal arteritis or polymyalgia rheumatica is suspected.
  \item \textbf{Complete blood count with differential:} To identify anemia (which elevates ESR), leukocytosis (infection/inflammation), thrombocytosis (reactive to inflammation), and abnormal differential (leukemia, lymphoma)
  \item \textbf{Peripheral blood smear:} When rouleaux are suspected from extremely elevated ESR (e.g., multiple myeloma)
  \item \textbf{Serum protein electrophoresis (SPEP) with immunofixation:} When multiple myeloma or Waldenstr\"{o}m macroglobulinemia is suspected (markedly elevated ESR with rouleaux on smear)
  \item \textbf{Serum immunoglobulins (IgG, IgA, IgM) and serum free light chains:} For monoclonal gammopathy evaluation
  \item \textbf{Rheumatoid factor and anti-cyclic citrullinated peptide (anti-CCP) antibodies:} When rheumatoid arthritis is suspected
  \item \textbf{Antinuclear antibody (ANA), anti-dsDNA, complement C3/C4:} When systemic lupus erythematosus or other connective tissue disease is suspected
  \item \textbf{Blood cultures and echocardiography:} When infective endocarditis is suspected
  \item \textbf{Imaging studies:} Chest X-ray, CT, MRI, or PET-CT for evaluation of suspected infection, malignancy, or inflammatory disease
  \item \textbf{Temporal artery biopsy:} The gold standard for diagnosis of temporal arteritis; should be performed as soon as possible (ideally within 1--2 weeks of starting corticosteroids, as histologic changes may persist for weeks after initiating therapy)
  \item \textbf{Fibrinogen level:} Directly measures the primary protein responsible for rouleaux formation and ESR elevation
  \item \textbf{Liver function tests and renal function tests:} To evaluate for underlying hepatic or renal disease affecting ESR
  \item \textbf{Tissue biopsy (lymph node, bone marrow, or organ-specific):} When malignancy is suspected
\end{itemize}

\section*{References}
\bibliographystyle{unsrtnat}
\bibliography{references}

\clearpage
